\documentclass[10pt,a4paper]{amsart}
\usepackage{amsmath,amssymb,amsthm,graphicx}
% Retained minimal preamble for the Stacks Project chapter more-algebra.tex (More on Algebra).
% The chapter file pinned-original-source.tex carries only \input{preamble} on its first line, so no
% preamble text lives inside the chapter itself.  The Stacks Project uses one single shared preamble for
% every chapter; the snapshot stages it identically as sources/stacks-chow/preamble.tex and
% sources/stacks-etale-cohomology/preamble.tex.  The block below is copied mechanically, line for line,
% from sources/stacks-chow/preamble.tex lines 14-18 (the Xy-pic setup) and 154-217 (the native theorem
% declarations together with the mathematical macros).  Nothing was added, removed or redefined; the
% wrapper class and the amsmath/amssymb/amsthm/graphicx base packages are supplied by the assembler.
% Chapter: more-algebra.tex -- More on Algebra.
% For commutative diagrams we use Xy-pic
\usepackage[all]{xy}

% We use 2cell for 2-commutative diagrams.
\xyoption{2cell}

% Theorem environments.
%
\theoremstyle{plain}
\newtheorem{theorem}[subsection]{Theorem}
\newtheorem{proposition}[subsection]{Proposition}
\newtheorem{lemma}[subsection]{Lemma}

\theoremstyle{definition}
\newtheorem{definition}[subsection]{Definition}
\newtheorem{example}[subsection]{Example}
\newtheorem{exercise}[subsection]{Exercise}
\newtheorem{situation}[subsection]{Situation}

\theoremstyle{remark}
\newtheorem{remark}[subsection]{Remark}
\newtheorem{remarks}[subsection]{Remarks}

\numberwithin{equation}{subsection}

% Macros
%
\def\lim{\mathop{\mathrm{lim}}\nolimits}
\def\colim{\mathop{\mathrm{colim}}\nolimits}
\def\Spec{\mathop{\mathrm{Spec}}}
\def\Hom{\mathop{\mathrm{Hom}}\nolimits}
\def\Ext{\mathop{\mathrm{Ext}}\nolimits}
\def\SheafHom{\mathop{\mathcal{H}\!\mathit{om}}\nolimits}
\def\SheafExt{\mathop{\mathcal{E}\!\mathit{xt}}\nolimits}
\def\Sch{\mathit{Sch}}
\def\Mor{\mathop{\mathrm{Mor}}\nolimits}
\def\Ob{\mathop{\mathrm{Ob}}\nolimits}
\def\Sh{\mathop{\mathit{Sh}}\nolimits}
\def\NL{\mathop{N\!L}\nolimits}
\def\CH{\mathop{\mathrm{CH}}\nolimits}
\def\proetale{{pro\text{-}\acute{e}tale}}
\def\etale{{\acute{e}tale}}
\def\QCoh{\mathit{QCoh}}
\def\Ker{\mathop{\mathrm{Ker}}}
\def\Im{\mathop{\mathrm{Im}}}
\def\Coker{\mathop{\mathrm{Coker}}}
\def\Coim{\mathop{\mathrm{Coim}}}

% Boxtimes
%
\DeclareMathSymbol{\boxtimes}{\mathbin}{AMSa}{"02}

%
% Macros for moduli stacks/spaces
%
\def\QCohstack{\mathcal{QC}\!\mathit{oh}}
\def\Cohstack{\mathcal{C}\!\mathit{oh}}
\def\Spacesstack{\mathcal{S}\!\mathit{paces}}
\def\Quotfunctor{\mathrm{Quot}}
\def\Hilbfunctor{\mathrm{Hilb}}
\def\Curvesstack{\mathcal{C}\!\mathit{urves}}
\def\Polarizedstack{\mathcal{P}\!\mathit{olarized}}
\def\Complexesstack{\mathcal{C}\!\mathit{omplexes}}
% \Pic is the operator that assigns to X its picard group, usage \Pic(X)
% \Picardstack_{X/B} denotes the Picard stack of X over B
% \Picardfunctor_{X/B} denotes the Picard functor of X over B
\def\Pic{\mathop{\mathrm{Pic}}\nolimits}
\def\Picardstack{\mathcal{P}\!\mathit{ic}}
\def\Picardfunctor{\mathrm{Pic}}
\def\Deformationcategory{\mathcal{D}\!\mathit{ef}}
\usepackage{cleveref}
\crefname{originallemma}{original lemma}{original lemmas}
\Crefname{originallemma}{Original lemma}{Original lemmas}
\makeatletter
\newlabel{more-algebra-lemma-cone-squared}{{062A}{source}{}{sourcepointer-1}{}}
\newlabel{more-algebra-lemma-cone-squared@cref}{{[originallemma][0][]062A}{[1][0][]source}{}{}{}}
\makeatother
\crefname{originallemma}{original lemma}{original lemmas}
\Crefname{originallemma}{Original lemma}{Original lemmas}
\makeatletter
\newlabel{equation-differential-cone}{{0627}{source}{}{sourcepointer-2}{}}
\newlabel{equation-differential-cone@cref}{{[originallemma][0][]0627}{[1][0][]source}{}{}{}}
\makeatother
\title{The cone on multiplication by a product is the cone of a map between the cones on the factors}
\author{The Stacks Project Authors}
\date{stacks-project:more-algebra.tex:local-snapshot-20261012; GNU FDL 1.2 or later}
\begin{document}
\maketitle
\setcounter{subsection}{1}
\noindent\textit{Editorial note.} All delivered mathematical text is the author's own, copied from the pinned Stacks Project chapter file more-algebra.tex; no statement or proof step was rewritten, reordered or removed.\par\medskip
\noindent\textit{Editorial note.} The chapter file pinned-original-source.tex carries only \textbackslash{}input\{preamble\} on its first line, so no preamble text lives inside the chapter. The retained file sources/stacks-more-algebra/retained-preamble.tex is a mechanical, line-for-line copy of the project's single shared preamble as staged in the snapshot (Xy-pic setup plus theorem declarations and mathematical macros). Nothing was added, removed or redefined, and the excerpt serves only as preamble, never as a span coordinate.\par\medskip
\noindent\textit{Editorial note.} Both counted spans start at the \textbackslash{}begin line of their environment, so the statement and the author's proof are delivered whole; surrounding prose of the section is deliberately outside the counted spans and is carried as standard prerequisites instead.\par\medskip
\noindent\textit{Editorial note.} The citation fields (source\_id, source\_version, authors, source\_url, license.name, source\_title) are identical to those of every other card this lane has produced from this chapter, so the batch selector treats them as one source.\par\medskip
\noindent\textit{Allowed standard prerequisites (editorial).}\begin{itemize}
\item Cone of a map of complexes and its differential, as recorded in the equation displayed earlier in the chapter. the author cites equation-differential-cone inside the proof at proof lines 6798-6801. Supplies the differential of the cone used in the special case.
\item Totalisation of a double complex and its compatibility with homotopies and cones. used at proof lines 6825-6841 to reduce to the special case. Carries the statement from the ring to an arbitrary complex.
\item Homotopy equivalence of chain complexes, verified here by exhibiting two maps whose composite is the identity on one side. used at proof lines 6805-6825. Supplies the criterion that the exhibited maps have to satisfy.
\end{itemize}\medskip
\ifdefined\hypertarget\hypertarget{sourcepointer-1}{}\fi
\noindent\textit{Source reference.} Original lemma 062A: focused-seeds/AGP/sources/stacks-snapshot-20261011/tags.txt, source lines 7789--7789. The lemma's own canonical tag, quoted so that readers can resolve the label inside the delivered text; it points at the counted statement itself and carries no proof content.\par\medskip
\ifdefined\hypertarget\hypertarget{sourcepointer-2}{}\fi
\noindent\textit{Source reference.} Original lemma 0627: focused-seeds/AGP/sources/stacks-snapshot-20261011/tags.txt, source lines 7786--7786. Navigation pointer to a statement the counted proof refers to but does not contain; nothing load-bearing is delegated to it.\par\medskip
\begin{lemma}
\label{lemma-cone-squared}
Let $R$ be a ring. Let $A_\bullet$ be a complex of $R$-modules.
Let $f, g \in R$. Let $C(f)_\bullet$ be the cone of
$f : A_\bullet \to A_\bullet$. Define similarly $C(g)_\bullet$ and
$C(fg)_\bullet$. Then $C(fg)_\bullet$ is homotopy equivalent to the
cone of a map
$$
C(f)_\bullet[1] \longrightarrow C(g)_\bullet
$$
\end{lemma}

\begin{proof}
We first prove this if $A_\bullet$ is the complex consisting of $R$ placed
in degree $0$. In this case the complex $C(f)_\bullet$ is the
complex
$$
\ldots \to 0 \to R \xrightarrow{f} R \to 0 \to \ldots
$$
with $R$ placed in (homological) degrees $1$ and $0$. The map
of complexes we use is
$$
\xymatrix{
0 \ar[r] \ar[d] &
0 \ar[r] \ar[d] &
R \ar[r]^f \ar[d]^1 &
R \ar[r] \ar[d] & 0 \ar[d] \\
0 \ar[r] & R \ar[r]^g & R \ar[r] & 0 \ar[r] & 0
}
$$
The cone of this is the chain complex consisting of $R^{\oplus 2}$ placed in
degrees $1$ and $0$ and differential (\ref{equation-differential-cone})
$$
\left(
\begin{matrix}
g & 1 \\
0 & -f
\end{matrix}
\right) :
R^{\oplus 2} \longrightarrow R^{\oplus 2}
$$
To see this chain complex is homotopic to
$C(fg)_\bullet$, i.e., to $R \xrightarrow{fg} R$,
consider the maps of complexes
$$
\xymatrix{
R \ar[d]_{(1, -g)} \ar[r]_{fg} & R \ar[d]^{(0, 1)} \\
R^{\oplus 2} \ar[r] & R^{\oplus 2}
}
\quad\quad
\xymatrix{
R^{\oplus 2} \ar[d]_{(1, 0)} \ar[r] & R^{\oplus 2} \ar[d]^{(f, 1)} \\
R \ar[r]^{fg} & R
}
$$
with obvious notation.
The composition of these two maps in one direction is the
identity on $C(fg)_\bullet$, but in the other direction
it isn't the identity. We omit writing out the required homotopy.

\medskip\noindent
To see the result holds in general, we use that we have a functor
$K_\bullet \mapsto \text{Tot}(A_\bullet \otimes_R K_\bullet)$
on the category of complexes which is compatible with homotopies
and cones. Then we write $C(f)_\bullet$ and $C(g)_\bullet$
as the total complex of the double complexes
$$
(R \xrightarrow{f} R) \otimes_R A_\bullet
\quad\text{and}\quad
(R \xrightarrow{g} R) \otimes_R A_\bullet
$$
and in this way we deduce the result from the special case discussed above.
Some details omitted.
\end{proof}

\end{document}
